% Teil II: durchgerechneter Weg, Teile a--d.

\subsection*{Gegeben und gesucht}
K\"orper auf schiefer Ebene (Winkel $\alpha$); Haftreibung
$F_R\le\mu_H F_N$, Gleitreibung $F_R=\mu_G F_N$. Gegeben für den Holzklotz:
$\alpha_H=27^\circ$, $\alpha_G=17^\circ$; f\"ur Autoreifen $\mu_H\approx0,8$
(trocken) bzw.\ 0,1 (eisig). Gesucht: (a) Zeigen $\mu_H=\tan\alpha_H$,
(b) experimentelle Bestimmung von $\mu_H,\mu_G$, (c) $\mu_H,\mu_G$ und
Beschleunigung bei $27^\circ$, (d) maximale Strassensteigung und die Rolle
von Masse/Auflagefläche.

\subsection*{Ansatz}
Kräftegleichgewicht senkrecht und parallel zur Ebene; Reibungsgrenze im
Moment vor dem Rutschen; gleichförmiges Gleiten mit $a=0$.
Beschleunigung nach dem Anstoss aus $a=g(\sin\alpha-\mu_G\cos\alpha)$.

\subsection*{Rechnung}
\textbf{a) Haftreibungswinkel.} Für den ruhenden Körper gilt parallel
zur Ebene $F_\parallel=F_R$; die Haftreibung erreicht gerade ihren
Höchstwert $\mu_H F_N$ im Moment des Rutschens:
\begin{equation}
  mg\sin\alpha_H=\mu_H\,mg\cos\alpha_H
  \;\Rightarrow\;
  \boxed{\;\mu_H=\tan\alpha_H\;}.
\end{equation}

\textbf{b) Experimentelle Bestimmung.} Man erhöht die Neigung langsam,
bis der Körper zu rutschen beginnt (Winkel $\alpha_H$), dann ist
$\mu_H=\tan\alpha_H$. Für $\mu_G$ stösst man den Körper an und findet
den Winkel $\alpha_G$, bei dem er gleichförmig (konstante Geschwindigkeit)
weitergleitet; dann ist $\mu_G=\tan\alpha_G$.

\textbf{c) Holzklotz.}
\begin{equation}
  \mu_H=\tan27^\circ\approx0{,}51,\qquad
  \mu_G=\tan17^\circ\approx0{,}31\;.
\end{equation}
Nach dem Anstoss bei $\alpha=27^\circ$ gilt mit der Gleitreibung
\begin{equation}
  a=g(\sin\alpha-\mu_G\cos\alpha)
  =9{,}81\,(\sin27^\circ-0{,}3057\cos27^\circ)
\end{equation}
\begin{equation}
  \sin27^\circ\approx0{,}4540,\ \cos27^\circ\approx0{,}8910,
  \;\Rightarrow\;
  a\approx9{,}81(0{,}4540-0{,}2724)\approx
  \boxed{\;\qty{1.78}{m/s^2}\;}.
\end{equation}
W\"are $\mu_G>\mu_H$ (mehr Gleit- als Haftreibung), so würde zwischen
$\alpha_H$ und $\alpha_G$ mit $\mu_H<\mu_G$ ein paradoxer Bereich
entstehen: der Körper beginnt beim Haftreibungswinkel zu rutschen, könnte
aber dann wegen der grösseren Gleitreibung ($a<0$ zwischen den beiden
Winkeln) nicht gleichförmig weiter; ein zunächst angefangenes Rutschen
kaeme wieder zum Stillstand -- unphysikalisch. Es muss stets
$\mu_H\ge\mu_G$ gelten.

\textbf{d) Maximale Strassensteigung.} Ein blockiertes (parkiertes) Rad
wird durch die Haftreibung gehalten; Rutschen setzt bei
$\tan\alpha=\mu_H$ ein:
\begin{equation}
  \alpha_{\max}=\arctan\mu_H: \quad
  \text{trocken } \alpha_{\max}=\arctan0{,}8\approx\qty{39}{\degree},
  \quad \text{eisig } \alpha_{\max}=\arctan0{,}1\approx\qty{5{,}7}{\degree}.
\end{equation}
Die Masse kürzt sich heraus ($mg$ auf beiden Seiten der
Reibungsgleichung); die Auflagefläche spielt im einfachen Modell
$F_R=\mu F_N$ keine Rolle.

\begin{ergebnis}
$\mu_H=\tan\alpha_H$, $\mu_G=\tan\alpha_G$; Klotz: $\mu_H\approx0{,}51$,
$\mu_G\approx0{,}31$, $a\approx\qty{1.78}{m/s^2}$ bei $27^\circ$.
Strassen: $\alpha_{\max}\approx\qty{39}{\degree}$ (trocken),
$\qty{5{,}7}{\degree}$ (Eis); Masse und Auflagefläche spielen im Modell
keine Rolle; $\mu_G>\mu_H$ w\"are unphysikalisch.
\end{ergebnis}

\begin{kontrolle}
\textbf{Einheiten:} $\tan$ als Verhältnis der Kraftkomponenten
dimensionslos; $a=g(\sin\alpha-\mu_G\cos\alpha)$ in \unit{\metre\per
\second\squared}. \checkmark

\textbf{Grenzfall:} Für $\mu_G=0$ (glatte Bahn) ist
$a=g\sin\alpha$, der freie Fall auf der Ebene. Für $\alpha=0$ (waagrecht)
wird $a=0$. \checkmark

\textbf{Plausibilitaet:} $\mu_H>\mu_G$ (0,51 > 0,31) erfuellt die
Forderung $\mu_H\ge\mu_G$; die Beschleunigung ist moderat, wie bei
flacher Reibung zu erwarten. \checkmark
\end{kontrolle}

\begin{merksatz}
Haftreibung hält bis $\mu_H=\tan\alpha_H$, Gleitreibung ist konstant
$\mu_G=\tan\alpha_G$; nach dem Anstoss beschleunigt der Körper mit
$a=g(\sin\alpha-\mu_G\cos\alpha)$. Masse und Fläche spielen im
Coulomb-Modell keine Rolle; es gilt stets $\mu_H\ge\mu_G$.
\end{merksatz}
